% ==============================================================================
% 07_forensics_memoire_reseau.tex — UE SIS589
% ==============================================================================
\chapter{Forensique Mémoire, Réseau et Cloud}
\label{chap:forensics-mem-net}

\epigraph{\itshape « La mémoire vive dit la vérité.
  Le disque peut mentir~; la RAM non. »}
         {--- Michael Cohen, auteur de Rekall}

\begin{boxobjectifs}
\begin{enumerate}
  \item Analyser un dump mémoire avec Volatility~3~: processus, connexions,
        DLL, injection mémoire.
  \item Détecter les techniques d'évasion en mémoire~: process hollowing,
        DLL injection, DKOM.
  \item Capturer et analyser le trafic réseau avec Wireshark et Zeek.
  \item Reconstituer une session réseau et extraire des fichiers transférés.
  \item Décrire les spécificités et défis de la forensique cloud (AWS,
        Azure, GCP).
\end{enumerate}
\end{boxobjectifs}

% ==============================================================================
\section{Forensique mémoire}
\label{sec:forensics-ram}
% ==============================================================================

\subsection{Pourquoi analyser la mémoire vive~?}

La RAM contient des informations qu'aucun autre artefact ne peut fournir~:
\begin{itemize}
  \item \textbf{Malwares fileless~:} exécutés entièrement en mémoire,
        sans fichier sur disque (T1059.001 + T1055). Invisibles au disque~;
        visibles en RAM.
  \item \textbf{Clés de chiffrement actives~:} présentes en mémoire
        pendant l'exécution~; disparaissent à l'extinction.
  \item \textbf{Processus injectés~:} code malveillant injecté dans un
        processus légitime (\texttt{svchost.exe}, \texttt{explorer.exe}).
  \item \textbf{Connexions réseau actives~:} l'état exact du réseau
        au moment de la saisie.
  \item \textbf{Mots de passe en clair~:} certains services (WDigest)
        maintiennent les credentials en clair en mémoire.
\end{itemize}

\subsection{Volatility~3~: framework d'analyse mémoire}

\begin{boxdef}
\textbf{Volatility~3~:} Framework open-source Python pour l'analyse
forensique de dumps mémoire Windows, Linux et macOS. Remplace Volatility~2
depuis 2020. Utilise des tables de symboles (ISF) spécifiques à chaque
version de noyau pour extraire les structures de données en mémoire.
\end{boxdef}

\begin{lstlisting}[style=StyleTerminal,
  caption={Commandes Volatility~3 essentielles pour l'investigation},
  label={lst:volatility3}]
# Toutes les commandes : python3 vol.py -f dump.mem windows.PLUGIN ou linux.PLUGIN

# ─── Identification du profil (système d'exploitation) ───────────────────────
python3 vol.py -f ram-dump-srv-web-01.raw windows.info
# Affiche : OS, version, build, architecture, date de dump

# ─── Liste des processus (pslist) ────────────────────────────────────────────
python3 vol.py -f ram.mem windows.pslist
# Colonnes : PID, PPID, ImageFileName, CreateTime, ExitTime

# pstree : vue arborescente (PPID → PID)
python3 vol.py -f ram.mem windows.pstree

# psscan : recherche directe dans la mémoire (détecte processus cachés)
python3 vol.py -f ram.mem windows.psscan
# Comparer pslist vs psscan : un processus dans psscan mais pas pslist
# = processus caché par un rootkit (DKOM)

# ─── Connexions réseau actives et historique ─────────────────────────────────
python3 vol.py -f ram.mem windows.netstat
# Affiche : PID, processus, protocole, src IP:port, dst IP:port, état

# ─── DLL chargées par processus ──────────────────────────────────────────────
python3 vol.py -f ram.mem windows.dlllist --pid 1234
# Toute DLL sans chemin ou dans %TEMP% est suspecte

# ─── Détecter les injections de code (code malveillant dans mémoire) ─────────
python3 vol.py -f ram.mem windows.malfind
# Recherche : pages mémoire exécutables avec des permissions anormales
# (READ|WRITE|EXECUTE sans fichier mappé)
# Résultat : dump hexadécimal + désassemblage des zones suspectes

# ─── Extraire les hachages de mots de passe (pour vérification) ──────────────
python3 vol.py -f ram.mem windows.hashdump
# NTLM hashes des comptes locaux (à croiser avec le SAM)

# ─── Lister les connexions réseau établies ───────────────────────────────────
python3 vol.py -f ram.mem windows.netscan
# Plus complet que netstat : inclut les connexions fermées récemment

# ─── Linux : équivalents ──────────────────────────────────────────────────────
python3 vol.py -f ram.mem linux.pslist
python3 vol.py -f ram.mem linux.pstree
python3 vol.py -f ram.mem linux.netstat
python3 vol.py -f ram.mem linux.bash  # historique bash en mémoire
\end{lstlisting}

\subsection{Détection des techniques d'évasion en mémoire}

\subsubsection{Process Injection (T1055)}

\begin{boxdef}
\textbf{Process Injection~:} Insertion de code malveillant dans l'espace
d'adressage d'un processus légitime pour masquer son exécution. Le processus
hôte (\texttt{svchost.exe}, \texttt{explorer.exe}) exécute le code de
l'attaquant avec ses propres privilèges.
\end{boxdef}

Indicateurs de détection avec Volatility~:
\begin{itemize}
  \item Pages mémoire \texttt{MEM\_PRIVATE + PAGE\_EXECUTE\_READWRITE}
        dans un processus normalement sans code privé (malfind).
  \item DLL injectée sans chemin fichier ou avec un nom proche d'une DLL
        légitime (\textit{DLL hijacking}).
  \item Threads distants créés dans un autre processus (CreateRemoteThread).
\end{itemize}

\subsubsection{Process Hollowing (T1055.012)}

\begin{boxdef}
\textbf{Process Hollowing~:} Un processus légitime est démarré en état
suspendu~; son code en mémoire est remplacé par du code malveillant~;
puis le processus est repris. Le gestionnaire de tâches affiche
\texttt{svchost.exe}~; le code exécuté est celui du malware.
\end{boxdef}

Indicateurs~: l'adresse de base de l'image en mémoire ne correspond
pas à l'exécutable sur disque (section \texttt{malfind}) ; checksum MZ/PE
divergent.

\subsubsection{DKOM (Direct Kernel Object Manipulation, T1014)}

\begin{boxdef}
\textbf{DKOM~:} Rootkit modifiant directement les structures du noyau
(liste doublement chaînée des EPROCESS) pour dissimuler un processus.
Le processus apparaît dans \texttt{psscan} (scan brute de la mémoire)
mais pas dans \texttt{pslist} (suit la liste officielle du noyau).
\end{boxdef}

\begin{lstlisting}[style=StyleTerminal,
  caption={Détection DKOM --- comparaison pslist vs psscan},
  label={lst:dkom-detection}]
# Extraire pslist et psscan, identifier les divergences
python3 vol.py -f ram.mem windows.pslist > pslist.txt
python3 vol.py -f ram.mem windows.psscan > psscan.txt

# Trouver les PIDs présents dans psscan mais absents de pslist
python3 - << 'EOF'
import re
with open('pslist.txt') as f:
    pids_list = set(re.findall(r'\b(\d{3,6})\b', f.read()))
with open('psscan.txt') as f:
    pids_scan = set(re.findall(r'\b(\d{3,6})\b', f.read()))
hidden = pids_scan - pids_list
print("Processus potentiellement cachés (DKOM) :", hidden)
EOF

# Si résultat non vide : processus caché = fort indicateur de rootkit
\end{lstlisting}

% ==============================================================================
\section{Forensique réseau}
\label{sec:forensics-net}
% ==============================================================================

\subsection{Capture du trafic réseau}

\begin{lstlisting}[style=StyleTerminal,
  caption={Capture réseau avec tcpdump et Wireshark},
  label={lst:capture-reseau}]
# ─── tcpdump : capture en ligne de commande ───────────────────────────────────
# Capturer tout le trafic vers/depuis srv-web-01 (10.0.1.10)
tcpdump -i em0 -w /mnt/evidence/capture-incident-20260412.pcap \
  host 10.0.1.10 \
  -C 100 \     # Rotation tous les 100 Mo
  -G 3600      # Rotation toutes les heures

# Capturer uniquement le trafic sortant suspect
tcpdump -i em0 -w capture.pcap \
  'src net 10.0.0.0/8 and not dst net 10.0.0.0/8 and not dst net 192.168.0.0/16'

# ─── tshark (Wireshark en ligne de commande) ─────────────────────────────────
# Afficher les connexions HTTP (sans SSL) depuis la capture
tshark -r capture.pcap -Y "http" \
  -T fields -e ip.src -e ip.dst -e http.request.method \
            -e http.request.full_uri -e http.response.code

# Extraire les fichiers transférés depuis le PCAP
tshark -r capture.pcap --export-objects http,/mnt/evidence/http_objects/
\end{lstlisting}

\subsection{Analyse avec Wireshark}

Filtres Wireshark essentiels pour l'investigation~:

\begin{lstlisting}[style=StyleCode,
  caption={Filtres Wireshark pour l'investigation forensique},
  label={lst:wireshark-filtres}]
# Trafic depuis/vers une IP suspecte
ip.addr == 203.0.113.42

# Connexions vers un port non standard (hors 80, 443, 22, 25)
tcp.port != 80 and tcp.port != 443 and tcp.port != 22 and tcp.dstport > 1024

# Requêtes DNS suspects (requêtes de type TXT = tunneling possible)
dns.qry.type == 16 or dns.qry.name contains "base64"

# Fichiers transférés en HTTP (POST avec body)
http.request.method == "POST" and http.content_length > 0

# Connexions TLS vers destinations non communes (JA3 hash)
ssl.handshake.type == 1   # ClientHello — vérifier JA3

# Détection de beacon C2 (connexions régulières vers une même IP)
# Utiliser Statistics > Conversations puis filtrer par interval régulier

# Reconstitution d'une session TCP complète
# Clic droit sur un paquet → Follow → TCP Stream
# → Affiche le dialogue complet en clair (si non chiffré)
\end{lstlisting}

\subsection{Zeek (anciennement Bro)~: analyse comportementale réseau}

\begin{boxdef}
\textbf{Zeek~:} Framework d'analyse de trafic réseau orienté sécurité.
Génère des journaux structurés (JSON) par protocole~: \texttt{conn.log}
(connexions), \texttt{dns.log}, \texttt{http.log}, \texttt{ssl.log},
\texttt{files.log} (fichiers extraits). Plus haut niveau que Wireshark~:
idéal pour la corrélation SIEM et la détection comportementale.
\end{boxdef}

\begin{lstlisting}[style=StyleTerminal,
  caption={Analyse Zeek sur une capture PCAP},
  label={lst:zeek-analysis}]
# Analyser un PCAP avec Zeek
zeek -r capture-incident-20260412.pcap local

# Fichiers générés : conn.log, dns.log, http.log, ssl.log, etc.

# ─── conn.log : toutes les connexions avec durée et volume ───────────────────
# Format : ts uid id.orig_h id.orig_p id.resp_h id.resp_p proto duration
# Identifier les connexions longues (beacon C2)
awk '$9 > 3600' conn.log | sort -k9 -rn | head -20

# Identifier les connexions répétées à intervalles réguliers (beacon)
awk '{print $4}' conn.log | sort | uniq -c | sort -rn | head -20

# ─── dns.log : requêtes DNS suspectes ────────────────────────────────────────
# Domaines avec entropie élevée (DGA)
awk '{print $10}' dns.log | sort -u | \
  awk '{
    l=length($0); e=0
    for(i=1;i<=l;i++) {c=substr($0,i,1); if(c~/[0-9]/) n++}
    if(n/l > 0.3) print $0, "HIGH_ENTROPY"
  }'

# ─── http.log : POST vers des URLs non standards ─────────────────────────────
awk '$15=="POST" && $7 !~ /\.(html|php|asp|js|css)$/' http.log

# ─── ssl.log : connexions TLS vers IPs non connues ───────────────────────────
awk '$11 !~ /amazon|google|microsoft|akamai/' ssl.log | head -20
\end{lstlisting}

\subsection{Reconstruction de sessions et extraction de fichiers}

\begin{lstlisting}[style=StyleTerminal,
  caption={Extraction de fichiers depuis un PCAP},
  label={lst:pcap-extraction}]
# ─── NetworkMiner : extraction automatique GUI ────────────────────────────────
# Ouvrir le PCAP dans NetworkMiner > onglet Files
# Extrait automatiquement : images, documents, binaires transférés

# ─── tshark : extraction par protocole ───────────────────────────────────────
# HTTP
tshark -r capture.pcap --export-objects http,/evidence/files_http/
# SMB
tshark -r capture.pcap --export-objects smb,/evidence/files_smb/
# FTP (données)
tshark -r capture.pcap --export-objects ftp-data,/evidence/files_ftp/

# ─── Zeek : extraction de fichiers avec hash ─────────────────────────────────
zeek -r capture.pcap /opt/zeek/share/zeek/policy/frameworks/files/extract-all-files.zeek
# Génère extract_files/ avec les fichiers + files.log avec SHA-256

# ─── Vérification des fichiers extraits via VirusTotal (CLI) ─────────────────
for f in /evidence/files_http/*; do
  hash=$(sha256sum "$f" | cut -d' ' -f1)
  echo "$hash  $(basename $f)"
done
# Soumettre les hashes manuellement ou via API VT
\end{lstlisting}

% ==============================================================================
\section{Forensique cloud}
\label{sec:forensics-cloud}
% ==============================================================================

\subsection{Spécificités et défis}

La forensique cloud présente des défis fondamentaux différents de la
forensique traditionnelle~:

\begin{description}
  \item[\textbf{Absence de contrôle physique~:}] L'investigateur n'a
        pas accès au matériel physique. Les images disques brutes sont
        inaccessibles (fournisseur cloud).
  \item[\textbf{Élasticité~:}] Les instances peuvent être créées et
        détruites en secondes. Un attaquant peut effacer ses traces en
        supprimant l'instance.
  \item[\textbf{Multi-tenant~:}] Les logs et métadonnées sont gérés par
        le CSP (\textit{Cloud Service Provider}).
  \item[\textbf{Multi-région~:}] Les données peuvent être réparties sur
        plusieurs juridictions.
\end{description}

\subsection{Sources forensiques cloud}

\begin{table}[H]
\centering\small
\caption{Sources de logs forensiques par fournisseur cloud}
\label{tab:cloud-logs}
\rowcolors{2}{TableauPair}{TableauImpair}
\begin{tabular}{p{2.5cm}p{4cm}p{7cm}}
\toprule
\tableauHeader{CSP} & \tableauHeader{Service} &
\tableauHeader{Données disponibles} \\
\midrule
\textbf{AWS} & CloudTrail &
API calls (qui a fait quoi, depuis quelle IP, quand) \\
\textbf{AWS} & VPC Flow Logs &
Trafic réseau (src/dst IP, ports, accept/reject) \\
\textbf{AWS} & S3 Access Logs &
Accès aux buckets (lecture, écriture, suppression) \\
\textbf{Azure} & Azure Monitor / Activity Log &
Actions sur les ressources Azure (CRUD) \\
\textbf{Azure} & Azure AD Sign-in Logs &
Authentifications, conditions d'accès conditionnel \\
\textbf{GCP} & Cloud Audit Logs &
Activité d'administration et d'accès aux données \\
\textbf{Commun} & Agent OS (Wazuh, OSSec) &
Logs système de l'instance VM (si déployé) \\
\bottomrule
\end{tabular}
\end{table}

\begin{lstlisting}[style=StyleTerminal,
  caption={Analyse des logs CloudTrail AWS --- investigation d'incident},
  label={lst:cloudtrail}]
# ─── AWS CLI : récupérer les événements CloudTrail ────────────────────────────
aws cloudtrail lookup-events \
  --start-time 2026-04-12T02:00:00 \
  --end-time   2026-04-12T06:00:00 \
  --max-results 50 \
  --query 'Events[?Username==`backup-user`]'

# ─── Filtrer les appels CreateUser et AttachUserPolicy (escalade) ─────────────
aws cloudtrail lookup-events \
  --lookup-attributes AttributeKey=EventName,AttributeValue=CreateUser \
  --start-time 2026-04-12T00:00:00

# ─── Analyser les logs CloudTrail JSON stockés en S3 ─────────────────────────
# Télécharger le log
aws s3 cp s3://lizbiz-cloudtrail-logs/AWSLogs/123456/CloudTrail/\
    ap-south-1/2026/04/12/ . --recursive

# Analyser avec jq
cat *.json.gz | zcat | \
  jq -r '.Records[] | select(.sourceIPAddress != "AWS Internal") |
  [.eventTime, .userIdentity.userName, .eventName,
   .sourceIPAddress, .errorCode] | @csv'
\end{lstlisting}

\subsection{Snapshot et acquisition d'instance cloud}

\begin{lstlisting}[style=StyleTerminal,
  caption={Acquisition forensique d'une instance EC2 compromise},
  label={lst:ec2-forensics}]
# ─── Étape 1 : Snapshot du volume EBS AVANT l'arrêt ─────────────────────────
aws ec2 create-snapshot \
  --volume-id vol-0a1b2c3d4e5f \
  --description "forensic-snapshot-incident-20260412" \
  --tag-specifications 'ResourceType=snapshot,Tags=[{Key=incident,Value=IR-2026-001}]'

# ─── Étape 2 : Créer un volume à partir du snapshot (en lecture) ─────────────
aws ec2 create-volume \
  --availability-zone ap-south-1a \
  --snapshot-id snap-0123456789abcdef0 \
  --volume-type gp3

# ─── Étape 3 : Attacher le volume à une instance d'analyse SIFT ──────────────
aws ec2 attach-volume \
  --volume-id vol-0newvolume \
  --instance-id i-0siftworkstation \
  --device /dev/sdf
# Le volume est attaché en lecture seule (équivalent write blocker)

# ─── Étape 4 : Monter et analyser ────────────────────────────────────────────
sudo mount -o ro /dev/sdf1 /mnt/forensics/
# → Appliquer les mêmes techniques que ch. 6 (registre, artefacts, timeline)
\end{lstlisting}

% ==============================================================================
\section*{Résumé}
% ==============================================================================

\begin{boxresume}
\begin{itemize}
  \item Forensique RAM~: Volatility~3. Commandes clés~: pslist, psscan,
        pstree, netstat/netscan, dlllist, malfind, hashdump.
  \item Détection d'évasion~: process injection (malfind),
        process hollowing (MZ/PE mismatch), DKOM (pslist~$\neq$~psscan).
  \item Forensique réseau~: tcpdump/Wireshark (PCAP), Zeek (logs structurés
        JSON), NetworkMiner (extraction fichiers), tshark --export-objects.
  \item Forensique cloud~: CloudTrail (AWS), Activity Log (Azure),
        Cloud Audit Logs (GCP). Snapshot EBS avant arrêt instance.
  \item Défis cloud~: pas d'accès physique, élasticité, multi-tenant,
        multi-juridictionnel.
\end{itemize}
\end{boxresume}

\begin{boxexo}
\textbf{Ex.~7.1 — Analyse mémoire.}
À partir du dump \texttt{ram-srv-web-01.raw} fourni (Lab~2), utilisez
Volatility~3 pour~: (a)~lister les processus actifs au moment de la
compromission~; (b)~identifier les connexions réseau établies~;
(c)~exécuter malfind et documenter les zones suspectes~;
(d)~déterminer si un processus est caché (DKOM).

\medskip\textbf{Ex.~7.2 — Analyse réseau.}
À partir du fichier \texttt{capture-incident.pcap} fourni, identifiez~:
(a)~l'IP de commande et contrôle (C2) utilisée~; (b)~le protocole
de communication~; (c)~les données exfiltrées~; (d)~la durée totale
de l'activité C2. Utilisez Wireshark + Zeek.

\medskip\textbf{Ex.~7.3 — CloudTrail.}
Les logs CloudTrail de LiZbiZ-SIS (compte AWS dev) indiquent une
action \texttt{CreateAccessKey} à 03:22~UTC le 12/04/2026 depuis
une IP TOR. Quelles sont les 4~prochaines actions d'investigation~?
Quels logs consulter~?
\end{boxexo}
